\section{Numerical evidence, corrections and open questions}\label{sec:limits}
Theorem~\ref{thm:main} identifies the \emph{exact good-prime recipe}.
It does not turn each stored approximation to that recipe into exact data.
The code supplement preserves the earlier correction package and its logs
unchanged. No archived research program, prime scan, certificate generator
or numerical $L$-function evaluator was executed for this consolidation.
The previously supplied standalone implementation rechecked only the
bounded cubic certificate and Sturm comparison in Section~\ref{sec:cubic};
its results are retained separately in the supplied supplement. The local continuations used only the small exact algebra checks recorded
in their supplements. No archived point-counting scan or numerical
$L$-function experiment has been rerun for this consolidation.

\subsection{Evidence retained from the earlier computations}
The carried record contains direct exact small-prime counts, an exact
finite modular-comparison certificate, rounded FFT-derived larger tables,
and high-precision evaluations using those tables. These are different
evidence classes. An exact consistency check on rounded input does not
certify how the input was obtained. The proof of the reciprocal template
also does not validate every historical value of its first two coefficients.

The historical $1296$ model imposed $P_2(T)=P_3(T)=1$.
The later $432$ twist in the retained N14/N15 record used
$P_2(T)=1$ and $P_3(T)=1-3T$. The blanket statement in version 2
that both experiments imposed $P_2=P_3=1$ was therefore too broad.
The new local proofs establish these factors and conductor products for
$W_4^-$ and $W_4^+$, respectively; they do not use the numerical fits.
The historical phrase ``plus-sign recipe'' refers to inert trace $S_p+1$,
which is the \emph{literal minus descent} $W_4^-$ of this manuscript.
For that $1296$ model the recorded functional-equation sign is $+1$
and the central value is approximately $2.2413552$. Those two statements
remain numerical observations, not consequences of the local proofs.
The candidate completion is
\[
 \Lambda(s)=N^{s/2}\Gamma_{\C}(s)^2\mathcal L(s),\qquad
 \Gamma_{\C}(s)=2(2\pi)^{-s}\Gamma(s).
\]
The associated theta-kernel Mellin transform used in the code has the
normalization $\widehat\Lambda=\Lambda/4$. Working precision and small
finite-sum defects do not supply an interval enclosure or an infinite-tail
bound. The historical gamma-factor shape agrees with the independent Hodge
calculation in Section~\ref{sec:scholl}. The conductor is now independently established by the local arguments.
The gamma-factor agreement and conductor calculation do not prove the
functional-equation sign or analytic continuation.

Likewise, a symmetrized evaluator
\[
 \sum_{n\le N}a_n\{G(s,n,t_0)+\epsilon G(6-s,n,1/t_0)\}
\]
has built-in symmetry at $t_0=1$. For $s=3$ and $\epsilon=-1$ it
vanishes for any coefficient sequence. This is not an independent
analytic-rank test. The historical N18 failure caused by comparing an
approximately fifty-digit reference against a seventy-digit requirement
is retained as a failure; other reported agreements are not promoted to
certified analytic digits.

\subsection{Two exact algebraic checks and their limits}
\begin{lemma}[Reciprocal-quartic modulus criterion]\label{lem:weil}
Let $P>0$, $c_1,c_2\in\R$,
$D=c_1^2-4(c_2-2P)$ and $B=16P-c_1^2-D$.
Every root of $X^4-c_1X^3+c_2X^2-Pc_1X+P^2$ has modulus
$\sqrt P$ if and only if
\[
 D\ge0,\qquad B\ge0,\qquad B^2\ge4c_1^2D.
\]
\end{lemma}
\begin{proof}
Put $y=X+P/X$. The two $y$-values are the roots of
$y^2-c_1y+c_2-2P$. The required condition is that they are real
and both lie in $[-2\sqrt P,2\sqrt P]$, equivalently
$D\ge0$ and $|c_1|+\sqrt D\le4\sqrt P$.
Squaring with the nonnegative signs retained gives the stated test.
\end{proof}
This repairs the earlier floor-square-root validator; it tests specified
coefficients. Good-prime purity of the exact geometric factors follows
independently from Corollary~\ref{cor:descent-good-prime-purity};
it does not certify any stored coefficient table.

If the same quartic factors as
$(X^2-aX+P)(X^2-bX+P)$, coefficient comparison gives
$D=(a-b)^2$. A nonsquare rational $D$ excludes this model with
$a,b\in\Q$. If $a,b$ are conjugate in a fixed quadratic field,
every nonzero $D$ has that field's square class. Neither this limited
obstruction nor empirical trace distributions establishes general
irreducibility, rules out every inner twist, or determines the Sato--Tate group.
The independent proof in Section~\ref{sec:irreducibility} now establishes
strong absolute irreducibility for $\ell\ne3$; it does not rely on this test.

\subsection{Remaining mathematical questions}
The bad-prime factors and conductor exponents have been determined in
Sections~\ref{sec:local-three}--\ref{sec:conductors}. Their auxiliary-prime
scope is recorded there; they are no longer open obligations of this draft.
The following questions remain separate.
\begin{enumerate}
\item The strong-irreducibility argument uses $\ell\ne3$. Its extension
at coefficient prime $3$, a full bad-prime Weil--Deligne comparison, and
precise Sato--Tate or inner-twist classifications are not established here.
\item Identify any automorphic representation with all required
hypotheses checked. The manuscript does not establish quaternion
multiplication or transfer an adjacent automorphy theorem.
\item Establish analytic continuation and the functional equation by a
theorem, including its sign. A root-number calculation requires its own
local epsilon-factor inputs; the conductor alone does not determine it.
\item Fix rational period structures before interpreting critical
values or determinant ratios arithmetically. The retained numerical
central values are not certified analytic evaluations.
\item For character orders $4,6$ in the older second-moment work,
supply their separate geometric inputs. Their finite comparison rows do
not follow from the cubic-only proof as all-prime theorems.
\end{enumerate}
These are limits of the present conclusions, not instructions to pursue
all of them before this manuscript can be independently reviewed.
The rank-four object is not the rank-two $H^1$ of an elliptic curve, even
after a Tate twist. No comparison controlling the Mordell--Weil rank or
BSD leading coefficient of a specified elliptic curve has been established
here. A future motivic special-value formulation would require its own
precise twist, arithmetic groups and comparison theorem.
